\begin{textsample}{POS Dim 6 – vem\_ed\_27 – Score 8.00 – vem\_ed\_27/t144.txt}  \label{ex:f6_pos_148}
Coordenador dos PCIs / Cesu de Norte a Sul da América Latina Em novembro e \textbf{dezembro} de 2024, Osvaldo Succi Junior, coordenador dos PCIs / Cesu, realizou visitas técnicas, conferências e \textbf{workshops} nos extremos Norte e Sul da América Latina: México e Chile.
No México, fez a conferência \textbf{magistral} “La Inteligência Artificial em entornos COIL”, no dia 27 de novembro, durante a 2ª Jornada COIL UNAM 2024, cujo tema foi “Aprender por el mundo sin salir de casa”.
A \textbf{apresentação} está disponível a partir do minuto 43 no link da \textbf{gravação} no Facebook: https://www.facebook.com/watch/live/?ef=watch\_permalink\&v=1282050929490379 Além da visita técnica à Universidad Nacional Autónoma de México (UNAM) e da conferência \textbf{magistral}, Succi Junior ofereceu o \textbf{workshop} “Consolidación institucional de um programa COIL”.
A \textbf{programação} incluiu, ainda, a palestra “\textbf{Una} introducción a los proyectos COIL” para um grupo de professores da Faculdade de Estudos Superiores Zaragoza e a \textbf{apresentação} “COIL en la formación universitaria” para docentes da Coordenação de Universidade Aberta e Educação Digital. continuação No Chile, o coordenador dos PCIs / Cesu ministrou um \textbf{workshop} sobre internacionalização para administradores da DUOC UC em 10 de \textbf{dezembro}; também uma palestra sobre Inteligência Artificial Generativa para professores da universidade \textbf{chilena} (no dia 11) e uma \textbf{apresentação} no 1º Seminário Internacional de Learn Chile.
Esse último evento, realizado em 12 de \textbf{dezembro} e transmitido ao vivo pelo YouTube, abordou o tema “Do local ao global: conectando a docência técnica profissional para um mundo globalizado”.
As principais instituições de ensino técnico profissional do Chile participaram do evento: DUOC UC, Inacap, Universidad Técnica Federico Santa María e Santo Tomás.
O objetivo do seminário foi refletir sobre desafios e oportunidades da formação técnica-profissional em contextos internacionais.
A rede Learn Chile conta com cerca de 20 instituições de ensino superior, incluindo universidades e institutos técnico-profissionais, e busca impulsionar a internacionalização da oferta acadêmica \textbf{chilena} e posicionar o país como destino educativo de excelência para estudantes internacionais.
A \textbf{gravação} do evento está disponível em: https://www.youtube.com/live/x6eTux-lIXI

% matched lemmas: apresentação, chileno, dezembro, gravação, magistral, programação, una, workshop
\end{textsample}
